\section{Organização do Documento}
\label{sec:organizacao}

O restante deste trabalho está organizado em seis capítulos.

O \textbf{Capítulo~\ref{cap:revisao}} apresenta o referencial teórico necessário para a leitura autossuficiente do trabalho: a estrutura do mercado de energia elétrica brasileiro e o papel do PLD, o problema econômico do comercializador de energia e a gestão de portfólio de contratos, os fundamentos da programação estocástica multiestágio e o método PDDE.

O \textbf{Capítulo~\ref{cap:revisao_lit}} conduz uma revisão crítica da literatura, examinando os trabalhos relacionados em três eixos: otimização de portfólio de contratos de energia no contexto brasileiro, métodos de solução para programação estocástica multiestágio aplicados a sistemas de energia, e modelagem estocástica do PLD. O capítulo conclui com o posicionamento do presente trabalho em relação às lacunas identificadas.

O \textbf{Capítulo~\ref{cap:modelagem}} apresenta a formulação matemática formal do problema: a notação adotada, o modelo DEQ com toda a árvore de cenários explícita e as restrições de não antecipatividade, e a decomposição via PDDE com o grafo de política e a recursão de Bellman.

O \textbf{Capítulo~\ref{cap:implementacao}} descreve a metodologia computacional: a arquitetura geral do sistema, o pipeline de geração de dados em Python (cenários de PLD via PAR(1), produção física estocástica, carteira legada e oportunidades de negociação), e as implementações do DEQ em JuMP e da PDDE em SDDP.jl.

O \textbf{Capítulo~\ref{cap:experimentos}} apresenta os experimentos computacionais: a configuração experimental, a análise de escalabilidade comparando DEQ e PDDE em instâncias de diferentes portes, a análise da política de contratação gerada por cada abordagem e a discussão dos resultados.

O \textbf{Capítulo~\ref{cap:conclusao}} sintetiza as conclusões do trabalho, discute as limitações da abordagem adotada e aponta direções para trabalhos futuros.
